% =============================================================================
% 03_3_hauptteil_web_frontend.tex — Domäne: Web-Frontend
% =============================================================================
% Dieser Abschnitt dokumentiert das Web-Frontend (Dashboard / UI) des
% Geotracker-Projekts. Gliederung gemäss Schulvorgabe:
%   1. Analyse und Design
%   2. Softwaredokumentation
%   3. Reflexion
% =============================================================================

\section{Web-Frontend}
\label{sec:frontend}

% Kurze Einführung in das Web-Frontend
% Welches Framework / welche Technologien werden eingesetzt?
% Was zeigt das Frontend dem Benutzer?
<<Kurze Einführung in das Web-Frontend (z.\,B.\ React, Vue.js, Leaflet.js für
Kartenansicht) und dessen Funktion im Gesamtsystem (Echtzeit-Positionsanzeige,
Verlaufsdarstellung).>>

% --- 1. Analyse und Design ---------------------------------------------------
\subsection{Analyse und Design}

\subsubsection{Anforderungsanalyse}
% Welche Anforderungen hat das Frontend zu erfüllen?
% (Benutzerfreundlichkeit, Responsivität, Echtzeit-Updates)
<<Anforderungen an das Web-Frontend auflisten (z.\,B.\ Kartenansicht mit
Echtzeit-Position, Verlaufslinie, Filterfunktionen, mobile Kompatibilität).>>

\subsubsection{GUI-Skizzen und Mockups}
% Füge Skizzen oder Mockups der wichtigsten Ansichten ein.
% Zeige, wie die Benutzeroberfläche aufgebaut ist.
% \begin{figure}[htbp]
%   \centering
%   \includegraphics[width=0.85\textwidth]{mockup_dashboard}
%   \caption{GUI-Mockup des Geotracker-Dashboards}
%   \label{fig:mockup_dashboard}
% \end{figure}
<<GUI-Skizzen / Wireframes für die Hauptansichten (Dashboard, Detailansicht,
Einstellungen) hier einfügen und erläutern.>>

\subsubsection{Komponentendiagramm}
% Skizziere die Komponentenstruktur der Frontend-Anwendung (UML oder informell).
<<Komponentendiagramm des Frontends beschreiben: Welche Komponenten gibt es?
Wie kommunizieren sie miteinander und mit der API?>>

% --- 2. Softwaredokumentation ------------------------------------------------
\subsection{Softwaredokumentation}

\subsubsection{Beschreibung und Konzepte}
% Erkläre die wichtigsten Komponenten, Zustandsverwaltung, Datenabruf.
% Wie werden Echtzeit-Daten abgerufen (WebSocket, SSE, Polling)?
<<Hauptkomponenten des Frontends beschreiben (z.\,B.\ Kartenkomponente,
Positionsmarker, Datenabruf via WebSocket/REST). Wichtige Codeausschnitte
als Listings einfügen.>>

% Beispiel:
% \begin{lstlisting}[caption={WebSocket-Verbindung aufbauen},
%                    label={lst:websocket}, language=JavaScript]
%   const ws = new WebSocket('wss://api.example.com/live');
%   ws.onmessage = (event) => {
%     const { lat, lng } = JSON.parse(event.data);
%     updateMarker(lat, lng);
%   };
% \end{lstlisting}

\subsubsection{Schwierigkeiten und Lösungsansätze}
% Welche Herausforderungen traten bei der Frontend-Entwicklung auf?
% (z. B. Cross-Origin-Requests, Kartenperformance, mobile Darstellung)
<<Technische Herausforderungen dokumentieren (z.\,B.\ CORS-Probleme, Performance
bei vielen GPS-Punkten auf der Karte, Responsive Design) und Lösungsstrategien
beschreiben.>>

\subsubsection{Testkonzept}
% Wie wurde das Frontend getestet?
% (Manuelle Tests, automatisierte UI-Tests, Cross-Browser-Tests)
<<Testkonzept für das Web-Frontend beschreiben. Welche Browser und Geräte wurden
getestet? Welche automatisierten Tests wurden eingesetzt (z.\,B.\ Jest,
Playwright)?>>

% --- 3. Reflexion ------------------------------------------------------------
\subsection{Reflexion}
% Bewerte das Frontend kritisch.
% Ist die Benutzeroberfläche intuitiv? Was würdest du verbessern?
<<Reflexion zur Frontend-Entwicklung: Stärken (z.\,B.\ übersichtliche Karte),
Schwächen (z.\,B.\ fehlende Offline-Unterstützung) und Verbesserungsvorschläge.>>
